% !PRINTED 565
\weberpage{565}
Wenn man in dem $(n-1)$fachen Integral $\xi_i$ durch
$\sqrt{1-\alpha^2}\,\xi_i$ ersetzt, so findet man dafür den Werth
$\sqrt{(1-\alpha^2)}^{\,n-1} K_{n-1}$, und es findet sich also die
Ungleichung

\[
K_n > 2\alpha \sqrt{(1-\alpha^2)}^{\,n-1} K_{n-1}.
\]

Nun wollen wir $\alpha$ so wählen, dass $2\alpha
\sqrt{(1-\alpha^2)}^{\,n-1}$ so gross als möglich wird, und dafür
findet man nach der gewöhnlichen Regel durch Nullsetzen des
Differentialquotienten

\[
\alpha = \frac{1}{\sqrt{n}}.
\]

Daraus

\[
K_n > 2\, \frac{\sqrt{(n-1)}^{\,n-1}}{\sqrt{n}^{\,n}}\, K_{n-1}
\]

\[
\left(\frac{\sqrt{n}}{2}\right)^{n} K_n > \left(\frac{\sqrt{n-1}}{2}\right)^{n-1} K_{n-1}.
\]

Setzt man also

\[
K_n = \left(\frac{2}{\sqrt{\Theta_n\, n}}\right)^{n},
\]

so ist

\weq{23}{\Theta_n^n < \Theta_{n-1}^{n-1},}

also $\Theta_n^n$ eine mit wachsenden $n$ abnehmende Zahl. Zugleich
ergiebt sich aus (16)

\weq{24}{M < \Theta_n\, n\, \sqrt[n]{D}.}

Für $n = 2$ ist $K_n$ gleich dem Flächeninhalt des Kreises mit dem
Radius 1, d. h. gleich der \gesperrt{Ludolfischen} Zahl $\pi$, und
daraus ergiebt sich

\[
\Theta_2 = \frac{2}{\pi} < 0{,}636,\qquad \Theta_2^2 < 0{,}404,
\]

% [sic: Weber prints references (22) and (23); the bound uses (24) and (23).]
und es gilt also nach (22) und (23) die für alle Werthe von $n$,
die grösser als 1 sind, gültige Grenzbestimmung:

\weq{25}{M < n \sqrt[n]{0{,}404\, D}\ \footnotemark.}

\footnotetext{Für specielle Anwendungen ist es nützlich, die Grenze
so genau als möglich zu kennen. Dazu dient eine genaue Berechnung
des Integrals $K_n$ nach seiner Definition (17), (18), die sich
leicht nach verschiedenen Methoden der Integralrechnung, z. B.
mittelst des discontinuirlichen Factors von \gesperrt{Dirichlet},
ausführen lässt:

\[
K_n = \frac{\Gamma\left(\frac12\right)^{n}}{\Gamma\left(1+\frac{n}{2}\right)},
\]
}
